\chapter{演化式事件驱动回测引擎}
\index{backtest}\index{event-driven system}\index{order}\index{fill}\index{portfolio}

\begin{introduction}
  \item 把 CSV、策略、订单、成交、组合和统计串成一条可解释数据流
  \item 用配置对象加入费用与滑点，并以独立账本守住旧行为
  \item 通过四个公开 seam 和独立 clean build 证明项目可重构、可复现
  \item 用正确性门控的端到端 benchmark 产生有边界的测量证据
\end{introduction}

\chaptermeta{第 5、8--16 章；能读类、模板、显式错误、CMake 目标、行为测试、benchmark、并发与数值边界。}
{Python 回测常把 DataFrame、策略函数和账户字典放在同一进程；C++ 版本把不可信文本、领域值、拥有对象和只读快照写成类型与调用边界，但不会因此自动变成真实交易平台。}
{最危险的结果不是崩溃，而是模型越界后仍打印一个看似精确的收益率。本章对混合代码估值明确拒绝，对性能数字先验收手算账本。}

\section{从行情分析器到事件循环}

第 5 章的行情分析器回答“这个 CSV 中某个代码发生了什么”；回测引擎还要回答“事件到达后，谁能产生意图、谁能改变账本、最终数字如何复算”。最终数据流是
\[
\text{CSV}\rightarrow\texttt{MarketEvent}\rightarrow\texttt{Strategy}
\rightarrow\texttt{OrderIntent}\rightarrow\texttt{Fill}
\rightarrow\texttt{Portfolio}\rightarrow\texttt{PerformanceSummary}.
\]
箭头不是装饰。CSV reader 是不可信文本边界；策略只观察事件与只读快照；撮合器把订单意图变成成交事实；只有 Portfolio 能应用成交；统计只读取权益曲线和成交集合。若策略直接改现金，策略测试就无法与账本测试分开，失败也难以定位。

最终源码保存在 \path{project/}，早期的 \path{code/ch05}、\path{code/ch08} 等快照仍从自己的 CMake 入口构建。最终项目没有反向引用这些目录，因此后续重构不会改变读者已经见过的教学状态。

\begin{pythonbridge}
Python 的 generator 或事件循环同样能分阶段处理数据；差异不在语法是否有 class，而在契约是否可被绕过。C++ 的 \texttt{const PortfolioSnapshot\&} 表达只读借用，值类型 \texttt{OrderIntent} 表达独立结果；Python 需要靠不可变对象、类型检查和团队约定达到相似效果。

\pythoncompare{generator 与 C++ 事件循环都能分阶段处理；领域事件应保持不同类型。}{Python 依靠对象引用；C++ 值事件独立拥有数据，引用快照只读借用。}{Python 靠运行期与类型工具；C++ 事件接口在编译期区分。}{Python 错字段常在路径执行时出现；C++ 错事件类型无法匹配接口。}
\end{pythonbridge}

\section{把配置视为领域输入}

旧引擎以报价立即成交且没有成本。第 17 章加入两个小配置值：每笔固定费用和以 basis points 表示的滑点。它们先进入 \texttt{ExecutionConfig}，再随 \texttt{BacktestConfig} 传给引擎；策略阈值和下单量仍属于策略对象，不混入撮合配置。

\begin{lstlisting}[caption={省略上下文：成本配置沿目标边界传递}]
struct ExecutionConfig final {
  double fixed_fee{};
  double slippage_bps{};
};

struct BacktestConfig final {
  double initial_cash{};
  ExecutionConfig execution;
};
\end{lstlisting}

买单执行价为 \(p(1+s/10000)\)，卖单为 \(p(1-s/10000)\)，每个 Fill 保存实际执行价与费用。撮合配置拒绝负数、非有限值以及不小于 10000 bps 的滑点，防止卖价变成非正数。Portfolio 应用成交时只做一次账本更新：
\[
\text{cash}_{new}=\text{cash}_{old}-q_{signed}p_{fill}-fee.
\]

独立账本是测试事实来源。报价 100、买 10 股、滑点 10 bps、费用 1.25 时，执行价是 100.10；从 10000 现金出发，现金变为 8997.75，以报价 100 估值的权益为 9997.75。Python 的 \texttt{dict} 配置也能传这些值，但字符串键拼错可能拖到运行路径才暴露；C++ 聚合类型让字段名、类型和默认值参加编译检查，却仍须在构造边界验证数值组合。

\section{四个 seam 保护演化}

测试不读取 Portfolio 的私有 map，也不检查某个辅助函数被调用几次。四个 seam 分别回答调用者真正关心的问题：

\begin{enumerate}
  \item \textbf{行情输入：}规范 CSV 产生带时间、代码、价格和数量的事件；非法行产生带行号错误。
  \item \textbf{策略决策：}低价且空仓时产生订单；高价或已有持仓时返回正常的空值。
  \item \textbf{成交与组合：}代码、数量和流动性满足契约才成交；现金、持仓与权益遵守手算账本。
  \item \textbf{端到端回测：}固定事件序列产生确定成交、最终快照和绩效摘要。
\end{enumerate}

零成本基线为 99 买 25 股、末价 103：现金 7525、权益 10100、总收益 1\%。成本 tracer bullet 使用同一序列、1 元费用和 10 bps 滑点：执行价 99.099、现金 7521.525、权益 10096.525、总收益 0.96525\%。首个事件后的权益为 9996.525，所以最大回撤是 0.03475\%，不是零。这个差异正说明为什么 expected 不能调用生产统计函数重新生成。

这里保留小规模 \texttt{double} 账本，是为了与前面章节的教学快照连续，并不把它升级为生产金额契约。若资产规模、币种或最小报价单位改变，必须先迁移到整数定点或领域金额类型，再重新定义费用、滑点、舍入、溢出和跨币种估值测试；仅增加有限性检查无法解决精度损失。

\begin{lstlisting}[language=bash,numbers=none]
ctest --test-dir build -R \
  "market_data_input|strategy_decision|execution_portfolio|end_to_end"
\end{lstlisting}

Python 的 pytest fixture 可以建立相同 seam；C++ 额外需要决定测试程序如何成为 CMake target。测试名描述“成本进入现金账本”，不绑定“调用了 match”这类实现步骤；容器从 \texttt{unordered\_map} 改成别的实现时，公开事实没有变化，测试就不应失败。

\section{多资产与缺失价格：先拒绝错误答案}

Portfolio 已能按代码保存多个仓位，但当前 \texttt{snapshot(symbol, price)} 只接收一个代码的一项价格。若事件序列混入 AAPL 与 MSFT，再用最后事件价格计算“总权益”，会得到语法合法、金融含义错误的数字。因此 \texttt{BacktestEngine::run} 目前要求一次运行只有一个代码；混合输入稳定失败：

\begin{lstlisting}[numbers=none]
model-boundary-error single-symbol backtest requires one symbol; define a complete price map before enabling multi-asset valuation
\end{lstlisting}

故意失败源 \path{project/failures/mixed_symbols.cpp} 不进入默认 \texttt{all}，CTest 先单独构建它，再验收状态 2 与上面的诊断。根因是估值契约缺少完整价格映射，与 vector 能否容纳两个代码无关。

真正的多资产 tracer bullet 至少要增加 \texttt{map<Symbol, Price>}，并在三种缺失价格政策中明确选择：拒绝本次估值；使用不晚于估值时刻的最近价格并报告 stale age；或把缺失仓位排除且把结果标记为不完整。直接填零会把未知价格伪装成资产归零，静默沿用未来价格则会引入 look-ahead bias。

\section{订单生命周期：当前只有即时终态}

本教学引擎的订单生命周期只有一条同步调用：策略产生 OrderIntent，撮合器立即返回一个完整 Fill 或 \texttt{nullopt}。没有排队订单、订单 ID、部分成交、撤单、有效期、交易所确认、重复消息和异步拒绝。这里的 \texttt{nullopt} 只表示“本事件没有成交”，不能回答为何拒绝，也不能作为生产审计记录。

若要演化，先画状态机，例如
\[
\text{created}\rightarrow\text{accepted}\rightarrow
\{\text{partially-filled},\text{filled},\text{cancelled},\text{rejected}\}.
\]
然后为每条合法边和非法反向边提供事实，给订单与事件稳定 ID，并让重复成交消息具有幂等规则。并发队列只负责传递这些状态，不能替代状态机。Python 对象也需要相同协议；async/await 只改变调度方式，不会自动给出幂等与审计语义。

\section{独立构建与可测量运行}

\path{project/CMakeLists.txt} 是最终快照入口。下面的命令不依赖根工程已经生成的目标；它从空目录构建 CLI、四个 seam、失败实验和 benchmark：

\begin{lstlisting}[language=bash,numbers=none]
cmake -S project -B build/ch17 -DCMAKE_BUILD_TYPE=Release
cmake --build build/ch17 -j
ctest --test-dir build/ch17 --output-on-failure
./build/ch17/quant_backtest project/data/sample.csv 1.0 10.0
\end{lstlisting}

成本配置输出必须逐字为：
\begin{lstlisting}[numbers=none]
backtest-ok events=3 fills=1 cash=7521.525 equity=10096.525 return=0.965250% max-drawdown=0.034750% volatility=0.000018 fee-per-fill=1.000 slippage-bps=10.000
\end{lstlisting}

端到端 benchmark 先生成 20 万个单代码事件，执行 2 次预热和 11 个样本；每个样本都必须保持一笔成交与权益 10100，之后才报告中位数和 IQR。完整可运行实现如下：

\lstinputlisting[caption={账本门控的端到端回测 benchmark}]{project/benchmarks/backtest_benchmark.cpp}

这条测量包含事件循环、策略、撮合、组合、权益曲线分配和统计，但不含 CSV、磁盘、Python 边界或多线程队列。Debug 数字不用于性能结论；Release 报告还要记录 CPU、编译器、优化配置、原始 11 个样本和系统负载。

\section{教学引擎与生产平台的分界}

当前项目有固定单币种现金、单代码运行、同步立即全额成交、固定费用与确定性滑点；没有交易日历、公司行动、复权、借券、保证金、风险门、容量曲线、延迟、数据版本和结果血缘。样例收益只能证明账本实现符合本章模型，不能证明策略可交易，更不能外推实盘收益。

合理演化顺序是：完整价格映射与缺失政策；订单状态机和拒绝原因；数据版本与交易日历；风险与容量；可回放的异步协议；最后才根据 profiler 证据调整布局或并发。每一步新增一条穿过输入、领域、输出与测试的 tracer bullet，并保持旧四个 seam 绿色。


\section{章末三层练习}

\subsection*{随堂验证汇总}

\begin{enumerate}
  \item 沿一条低价事件写出每个阶段的输入类型、输出类型及所有者；任何一步若同时返回“事件和错误字符串”，说明边界还没有说清。
  \item 把费用改为 2.00，只改一项输入，先把期望现金写成 8997.00，再运行成交/组合 seam。若通过“放宽容差”让错误账本变绿，修复方向已经错了。
  \item 先把成本账本的期望权益故意改错 1 元，确认测试因事实不符而红，再恢复独立值。
  \item 为三种政策各写一个输入、可观察输出和风险，暂时不要修改引擎。
  \item 解释 \texttt{filled -> accepted} 为什么非法，以及部分成交后撤单时剩余数量和已成交数量各由谁拥有。
  \item 只把事件数加倍，检查账本仍相同，再比较中位数；一次运行变快不能证明复杂度改善。
\end{enumerate}

复现层重放三事件账本；修改层加入一个费用或失败输入并保持 seam 可测；开放层选择多资产估值、取消状态机或数据恢复之一，只交付最小纵向闭环。

输出任务：从空目录配置 \path{project/}，构建并运行全部 CTest；提交默认与成本配置的两条精确摘要、四个 seam 的通过记录、混合代码状态 2 诊断，以及 Release benchmark 的环境、11 个原始样本、中位数、IQR 和账本门结果。另附一页边界说明，覆盖费用、滑点、多资产、缺失价格和订单生命周期。

自测：
\begin{enumerate}
  \item 为什么 MarketEvent、OrderIntent 与 Fill 不能合成一个“万能行”类型？
  \item 10 bps 买入滑点如何从报价 99 得到执行价 99.099？
  \item 为什么成本路径会在第一条事件产生非零最大回撤？
  \item 四个公开 seam 各自保护什么行为，为什么不检查私有 map？
  \item 当前多资产输入的真正缺失契约是什么？
  \item \texttt{nullopt} 成交结果为什么不能替代生产订单状态机？
  \item clean build 比在旧目录再次运行 build 多证明了什么？
  \item benchmark 的账本门、预热、重复样本和局限分别解决什么问题？
\end{enumerate}

\section{本章小结}

最终项目的可信度来自可复算边界，而非文件数。输入先验证，策略只产生日志化意图，成交携带真实成本，组合独占账本修改，统计读取已定义的权益曲线；四个 seam、失败实验、clean build 和账本门控 benchmark 共同约束演化。明确拒绝尚未定义的多资产估值，比输出一个漂亮但错误的数字更接近生产工程。
